\section{Preliminaries}\label{sec:prel}
%\todo{Need to define pathwidth somewhere}
We use the following basic notation. For a positive integer $\ell$, we use $[\ell]$ to denote the set $\{1, 2, \dots, \ell\}$. For a graph $G$, we use $V(G)$ to denote the set of vertices of $G$, and $E(G)$ to denote the set of edges of $G$. For a graph $G$ and a subset $X \subseteq V(G)$, we write $G[X]$ to denote the subgraph of $G$ induced by $X$, i.e., a graph with vertices restricted to $X$ and all edges between vertices of $X$ preserved.

By $\sym_k$ we denote the group of all permutations on the set $[k]$. By $\bar\sym_k$ we denote the subset of $\sym_k$ of permutations that contain exactly one cycle.
By $\sym_{2k}^2$ we denote the subset of $\sym_{2k}$ of permutations that contain exactly $k$ cycles, each of size $2$ and by $\sym_{2k}^k$ we denote the subset of $\sym_{2k}$ of permutations that contain exactly $2$ cycles, each of size $k$.

A \emph{cycle type} of a permutation is defined as the multiset of the sizes of all its cycles.
A\emph{conjugation} of a permutation $\rho \in \sym_k$ by a permutation $\pi \in \sym_k$ is defined as $\pi^{-1} \rho \pi$, i.e. the composition of the inverse of $\pi$, $\rho$ and $\pi$. It is well known that permutation conjugation preserves its cycle type, and hence $\bar\sym_k, \sym_{2k}^2, \sym_{2k}^k$ are closed under taking conjugations.
We will need the following group-theoretic properties of permutations and their conjugations. All of those are either well known or easy to obtain with some basic calculations done on subgroups.

\begin{prop}\label{prop:sym_k}
    The following equalities hold for any $k \geq 1$ and for any maximal set $\mathcal{Z}_k \subseteq \sym_k$ of permutations of the same cycle type.
    \begin{enumerate}
        \item $|\bar\sym_k| = (k - 1)!$, $|\sym_{2k}^2| = (2k - 1)!!$, $|\sym_{2k}^k| = (2k - 1)! / k$,
        \item $\{ \pi^{-1}\rho\pi \mid \pi \in \sym_k \} = \mathcal{Z}_k$ for every $\rho \in \mathcal{Z}_k$,
        \item $|\{ \pi \in \sym_k \mid \pi^{-1} \rho_1 \pi = \rho_2 \}| = |\sym_k| / |\mathcal{Z}_k|$ for every $\rho_1, \rho_2 \in \mathcal{Z}_k$,
        \item $|\{ \rho_2 \in \sym_{2k}^2 \mid \rho_2 \rho_1 \in \sym_{2k}^k \}| = (2k - 2)!!$ for every $\rho_1 \in \sym_{2k}^2$.
    \end{enumerate}
\end{prop}

\paragraph{Pathwidth.}
Let $G$ be a graph (undirected or directed). A \emph{path decomposition} of $G$ is a sequence of subsets $\beta_1, \beta_2, \dots, \beta_\ell \subseteq V(G)$ called \emph{bags} such that:
\begin{itemize}
    \item for every edge $(u,v) \in E(G)$, there is some $i \in [\ell]$ such that $u,v \in \beta_i$, and
    \item for every vertex $v \in V(G)$, there exist $1 \leq i_1 \leq i_2 \leq \ell$ such that $\{ i \mid v \in \beta_i \} = \{i_1, i_1 + 1, \dots, i_2 - 1, i_2\}$, i.e., the bags containing $v$ form a connected subinterval of the sequence of all bags.
\end{itemize}
The \emph{width} of a path decomposition $(\beta_i)_{i=1}^\ell$ is defined as $\max_{i \in [\ell]} |\beta_i| - 1$. The \emph{pathwidth} of $G$ is defined as the minimum possible width of a path decomposition of $G$.


\subsection{Tropical polynomials}
We treat all polynomials as defined over the tropical $(\max, +)$-semiring, i.e., $(f \cdot g)(x) = f(x) + g(x)$ and $(f + g)(x) = \max(f(x), g(x))$ for any two polynomials $f, g$.
%We assume that all polynomials are over the same set of variables which is implicitly defined.

\paragraph{Independent Set Polynomial.} If we treat the weights $x_v$ as indeterminates we can view $IS_G$ as a polynomial in the (max,+) semiring by replacing the max operation by addition and the $+$ operation by multiplication:
\begin{equation}\label{eq:mwispoly}
    IS_G := \sum_{I \in \mathcal{I}(G)} \left(\prod_{v \in I} x_v\right),
\end{equation}
where we remind the reader that $\mathcal{I}(G)$ denotes the family of all independent sets of $G$.

\paragraph{Traveling Salesperson Problem polynomial.} Similarly to the Independent Set polynomial $IS_G$, we define the directed TSP polynomial of $G$ as
$$
DTSP_G := \sum_{u_0, \dots, u_N \in \mathcal{H}(G)}
\left( \prod_{i \in [N]} x_{u_{i - 1}, u_i} \right),
$$
where we remind the reader that $\mathcal{H}(G)$ denotes the family of all sequences of $N + 1$ vertices $(u_i)_{i =0}^N$ such that $u_0 = u_N$ and $u_0 \to u_1 \to \dots \to u_{N - 1} \to u_N$ is a directed Hamiltonian cycle in $G$.
Naturally, we also view $TSP_G$ similarly as a polynomial by identifying variables $x_{u,v}$ and $x_{v, u}$ for all $uv \in E(G)$.

\paragraph{Directed Spanning Tree polynomial.} We also define a directed spanning tree polynomial of $G$ as
$$
DST_G := \sum_{p \in \mathcal{T}(G)}
\left( \prod_{v \in V(G) : p(v) \neq v} x_{p(v), v} \right)
$$
where we remind the reader that $\mathcal{T}(G)$ denotes the family of all functions $p : V(G) \to V(G)$ such that for exactly one vertex $v \in V(G)$, $p(v) = v$, and edges $(p(u), u)$ over all other vertices $u$ form an out-tree rooted at $v$.
\newline

\noindent
The \emph{support} of a monomial $m = x_1^{a_1} \dots x_\ell^{a_\ell}$ is the set of variables $\{x_1, \dots, x_\ell\}$ and is denoted as $\sup(m)$. The \emph{support of a polynomial} $p = m_1 + \dots + m_\ell$ is the union of supports $\sup(m_i)$ over all $i \in [\ell]$.
We let $\sup(p)$ denote the support of polynomial $p$.
For a monomial $m$ and a polynomial $p$, we write $m \in p$ if $p$ treated as a formal expression is of the form $m + q$ for some polynomial $q$. We say that a polynomial $p$ is \emph{homogeneous} if and only if for some $d \in \mathbb{N}$, the degree of all monomials $m \in p$ is $d$. For a polynomial $p$, we write $|p|$ to denote the number of monomials $m$ such that $m \in p$.

A \emph{valuation} is a function which maps variables to $\mathbb{R}$.
A \emph{characteristic valuation $\chi_m$} of a monomial $m$ is a valuation such that $\chi_m(x) = 1$ if $x \in \sup(m)$ and $\chi_m(x) = -1$ otherwise. As~stated before, we evaluate polynomials in a $(\max, +)$-semiring, i.e., given a valuation $v$, the~monomial $x$ evaluates to $v(x)$, $(f + g)(v)$ evaluates to $\max(f(v), g(v))$ and $(f \cdot g)(v)$ evaluates to $f(v) + g(v)$.

For two polynomials $p, q$, we write $p \subseteq q$ if for every monomial $m \in p$, we have $m \in q$. We write $p \simeq q$ iff $p \subseteq q$ and $q \subseteq p$. In particular, $p \simeq q$ doesn't imply that $p = q$ (take e.g. $p = x$ and $q = x + x$).
The definition of $\simeq$ is motivated by the following observation.

\begin{prop}\label{prop:poly_simeq}
    For any two multilinear polynomials $f, g$, we have $f \simeq g$ if and only if for any valuation $v$, we have $f(v) = g(v)$.
\end{prop}

\begin{proof}
    The implication from $f \simeq g$ to $f(v) = g(v)$ is immediate. In the other direction, assume w.l.o.g. that $f \subseteq g$ does not hold, and hence there is some monomial $m \in f$ such that $m \not\in g$.
    It is easy to see that $f(\chi_m) = |\sup(m)|$ and $g(\chi_m) < |\sup(m)|$, which is a contradiction.
\end{proof}


\subsection{Tropical circuits}

A \emph{tropical circuit} is a directed acyclic graph in which every vertex (called \emph{node}) is of in-degree 0 or 2. Vertices with in-degree 0 are labeled with either a variable, in which case they are called \emph{input nodes}, or with a constant $0$, in which case they are called \emph{constant nodes}. The nodes with in-degree $2$ are called \emph{operation nodes} and are labeled with a binary operator, either $+$ or $\max$. There is one node, designated an \emph{output node}.

Evaluation of a tropical circuit given a valuation $v$ is defined the following way. The nodes are processed according to the topological order of the graph. Constant nodes evaluate to value $0$. Input nodes labeled $x$ evaluate to $v(x)$. Operation nodes labeled with, respectively, $+$ and $\max$ evaluate to, respectively, sum and maximum of the values of its two predecessors. The output of the evaluation of the whole circuit is the value obtained at the output node.

We say that a tropical circuit $\Gamma$ \emph{calculates some polynomial $p$}, if the evaluation of $\Gamma$ on $v$ is equal to $p(v)$ for every valuation of variables $v$. It is easy to see that semantically, evaluating a tropical circuit computes exactly some polynomial: input node labeled $x$ computes $x$, the $\max$ nodes compute the sum of two polynomials and the $+$ nodes compute the product of two polynomials.


By \Cref{prop:poly_simeq} and the definitions, we immediately get the following.

\begin{prop}\label{prop:tc_semantics}
    Let $f, g$ be two multilinear polynomials and let $\Gamma$ be a tropical circuit calculating~$f$. Then $\Gamma$ calculates $g$ iff $f \simeq g$.
\end{prop}
Intuitively, this claim states that a tropical circuit calculates some polynomial iff the set of its monomials is exactly the set of all monomials which appear during the evaluation of said circuit.

The main combinatorial ingredient regarding tropical circuits we will use is the following decomposition lemma. For the proof, see, e.g., \cite[Lemma 3.4]{Jukna2023}.
\begin{lemma}\label{lem:tc_decomposition}
    Let $f$ be a homogeneous polynomial calculated by a tropical circuit of size $\tau$ and let $X$ denote an arbitrary subset of $\sup(f)$. Then $f$ can be written as
    $$
        f \simeq \sum_{i \in [\tau]} g_i \cdot h_i,
    $$
    where $|\sup(g_i) \cap X|, |\sup(h_i) \cap X| \leq \frac{2}{3}|X|$ for each $i \in [\tau]$.
\end{lemma}



